\section{Report Structure}
\label{sec:intro_structure}

The remainder of this research monograph is organized as follows:
\begin{itemize}
    \item \textbf{Chapter 2 (Theoretical Foundations and Literature Review):} Reviews Probabilistic Graphical Models, Structural Causal Models, surveys DAG structure learning algorithms (PC, FCI, Hill Climb, GES, NOTEARS) and parameter learning methods (MLE, Bayesian BDeu, EM, Neural CPDs), and reviews semi-parametric estimators and educational data mining literature.
    \item \textbf{Chapter 3 (Methodology and System Architecture):} Formalizes the 7-node Primary Domain DAG, establishes the four-pillar DAG validity constraints, specifies the $2 \times 2$ factorial model space ($B_0, H_1, H_2, H_3$), details the Deep Neural CPD architecture and Temporal Forbidden Mask, and specifies the DoWhy refutation suite.
    \item \textbf{Chapter 4 (Empirical Evaluation and Core Benchmark Results):} Presents Table A (held-out predictive performance), Table B (causal interventions and ATE isolating spurious confounding), and Table C (adversarial DoWhy refutations and computational runtime), accompanied by ablation diagnostics.
    \item \textbf{Chapter 5 (Advanced Causal Extensions and Policy Optimization):} Explores continuous GPA and retention estimands, causal mediation path decompositions ($NDE$ vs $NIE$), heterogeneous treatment effects (CATE), Cinelli-Hazlett sensitivity contours ($RV_{q=1}$), budget-constrained policy optimization (+164 graduates), and large-scale Kaggle S4E6 out-of-distribution transferability.
    \item \textbf{Chapter 6 (Conclusions, Limitations, and Future Directions):} Synthesizes the core theoretical and empirical findings, verifies research hypotheses, articulates actionable institutional policy directives, assesses threats to validity, and outlines future research directions.
\end{itemize}
